\documentclass[fleqn]{article}
\usepackage{amsmath}
\usepackage{amssymb}
\usepackage{array}
\usepackage[left=2cm,right=2cm,top=1.5cm,bottom=1.5cm]{geometry}
\setlength{\mathindent}{1.5em}

\newcommand{\Rf}{R\bigl(B_i^{\,j}\bigr)}

\begin{document}

\begin{center}
{\LARGE \textbf{Radix-2}}\\[4pt]
{\large DP8 (8$\times$4) --- one digit per bit of $B$}
\end{center}

\vspace{1.5em}

\noindent\textbf{Step 1 --- the dot product}
\[
\sum_{i=0}^{7} A_i \, B_i
\]

\vspace{0.5em}

\noindent\textbf{Step 2 --- expand $B$ into its bits and exchange the sums}

\noindent Each $B_i$ is a 4-bit value read as signed: the top bit counts negative,
$B_i = \sum_{j=0}^{3} w_j \, B_i^{\,j}$ with $w = (1, 2, 4, -8)$, so
\[
\sum_{i=0}^{7} A_i \, B_i
  \;=\; \sum_{i=0}^{7} \Biggl[\, \sum_{j=0}^{3} w_j \, B_i^{\,j} \Biggr] A_i
  \;=\; \sum_{j=0}^{3} w_j \sum_{i=0}^{7} B_i^{\,j} \, A_i
\]

\noindent When $B_i$ is a lower slice of a wider $B$ the top bit is read unsigned, $w_3 = +8$.

\vspace{0.5em}

\noindent\textbf{Step 3 --- the partial product}
\[
\sum_{i=0}^{7} A_i \, B_i
  \;=\; \sum_{j=0}^{3} w_j \sum_{i=0}^{7}
        \underbrace{\Bigl[\, B_i^{\,j} \, A_i \Bigr]}_{\textstyle \Rf}
\]

\noindent $\Rf$ depends on one bit and takes one of two values; the sign never enters
the digit, it is the weight of the top plane, so exactly one row is subtracted:

\vspace{0.5em}

\begin{center}
\renewcommand{\arraystretch}{1.2}
\begin{tabular}{c | c | l}
$B_i^{\,j}$ & $d_i^{\,j}$ & \multicolumn{1}{c}{$R$} \\
\hline
$0$ & $0$ & $0$ \\
$1$ & $1$ & $A_i$ \\
\end{tabular}
\end{center}

\vspace{0.8em}

\noindent\textbf{Step 4 --- split the work}

\noindent There is nothing to encode: the digit is the bit itself, and the dispatcher sends
it to the PE as the control signal of the selection; only the selection stays inside:
\[
\underbrace{\bigl\{\, B_i^{\,j} \,\bigr\} \;\longrightarrow\; d_i^{\,j}}_{\textstyle \text{dispatcher: encoding}}
\qquad\longrightarrow\qquad
\underbrace{\sum_{j=0}^{3} w_j \sum_{i=0}^{7} d_i^{\,j} \, A_i}_{\textstyle \text{PE: selection}}
\]

\noindent Each $d_i^{\,j}$ travels as its bit, which the cell turns into one of $\{0, A_i\}$:
4 control bits per lane.

\vspace{0.8em}

\noindent $B$ is the narrow operand: 4 digits $\times$ 8 lanes = 32 selections of
8 bits per DP8, one multiple ($A_i$) and no encoder, at twice the rows of every
radix-4 scheme.

\end{document}
